\SourceLeaf{432}wenn der Nordpol des Magnetstabes auf der Westseite, und
\[
T=(1-\frac{n-13{,}65c}{16910})\mathfrak{T}
\]
wenn er auf der Ostseite sich befindet.

Übrigens bedarf es keiner Erinnerung, dass die Berücksichtigung der Temperatur bei den Terminsbeobachtungen füglich ganz unterbleiben kann, so lange man nur darauf ausgeht, die Gestaltung der einzelnen in kurzen Zeitfristen wechselnden Anomalien zu erkennen.

\Article{14}
Wie bei der Transversalstellung des Bifilarmagnetometers die Veränderungen der Intensität in ihrer ganzen Stärke, die der Declination hingegen kaum merklich den Stand afficiren, so haben gerade umgekehrt auf die Schwingungsdauer die letztern Veränderungen den bedeutendsten, die erstern hingegen nur einen äusserst geringen Einfluss. In sofern $p,q$ und die Abweichung des Elements $S$ von dem Normalwerthe nur klein sind, wird ohne erheblichen Fehler anstatt der Formel (2) Art. 11 gesetzt werden können
\[
\frac{t^0t^0}{tt}=1\mp q\operatorname{tang}\varphi\quad\text{oder auch}\quad t=t^0(1\pm\tfrac12q\operatorname{tang}\varphi)
\]
wenn $q$ in Theilen des Halbmessers, und folglich
\[
t=t^0\pm\frac{t^0\operatorname{tang}\varphi}{412530}.q
\]
wenn es in Bogensecunden ausgedrückt ist. Aus den Resultaten des oben berechneten Beispiels folgt $t^0=55{,}871$ Zeitsecunden, wonach also in Bogensecunden
\[
q=\pm(t-55{,}871)4172''8
\]
wird. Die ganz scharfe Transformation der Formel (2) zur Berechnung von $q$ ist folgende
\[
\pm\operatorname{tang}q=\frac{tt-st^0t^0\cos p}{tt\operatorname{tang}\varphi\pm st^0t^0\sin p}
\]
Übrigens bedarf es keiner Erinnerung, dass auf diese Weise durch Schwin\-gungsbeobachtungen nicht der für einen bestimmten Augenblick gültige Werth von $q$, sondern nur der Mittelwerth für die Dauer jener Beobachtungen bestimmt werden kann.

\clearpage
\fancyhead[L]{\small MAGNETISCHE WIRKUNG IN DER FERNE.}
\SourceLeaf{433}\pdfbookmark[0]{Vorschriften zur Bestimmung der magnetischen Wirkung welche ein Magnetstab in der Ferne ausübt.}{fernwirkung-title}
\begin{center}
{\Large\bfseries VORSCHRIFTEN ZUR BESTIMMUNG\par\medskip DER MAGNETISCHEN WIRKUNG\par\medskip WELCHE EIN MAGNETSTAB\par\medskip IN DER FERNE AUSÜBT.}
\par\bigskip
Resultate aus den Beobachtungen des magnetischen Vereins. 1840. II.
\end{center}
\renewcommand{\Article}[1]{\par\medskip\Needspace{4\baselineskip}\pdfbookmark[1]{#1.}{fernwirkung-art-#1}\centerline{\textbf{#1.}}\nobreak\smallskip}%
Wenn man mehrere magnetische Apparate zugleich aufgestellt hat, dürfen die gegenseitigen Einwirkungen nicht unbeachtet bleiben. Die verschiedenen Apparate in so grossen Entfernungen von einander aufzustellen, dass diese Einwirkungen unbesehens für ganz unmerklich geachtet werden können, ist ein nicht überall anwendbares, und jedenfalls mit der Aufopferung mancher sonstigen Vortheile und Bequemlichkeiten verknüpftes Auskunftsmittel. Kann man aber die Wirkungen, welche ein Apparat an dem Platze eines andern ausübt, durch Rechnung mit Schärfe bestimmen, und also von den am zweiten Apparate gemachten Beobachtungen abtrennen, so behält man die vollkommenste Freiheit, bei der Wahl der Aufstellungsplätze jeder andern Rücksicht ihr Recht widerfahren zu lassen, und die Entwickelung der zu diesen Rechnungen dienenden Formeln scheint daher hier einen Platz wohl zu verdienen.

\Article{1}
Die Lage des Punktes, für welchen die Wirkung eines Magnetstabes berechnet werden soll, werde durch drei rechtwinklige Coordinaten, $x,y,z$ bestimmt, deren Anfang wir hier in den Mittelpunkt des Magnetstabes selbst setzen; um die Vorstellungen zu fixiren, nehmen wir an, dass die beiden ersten Coordinatenaxen horizontal sind, und zwar die erste im wahren Meridiane, die dritte
\SourceLeaf{434}also vertical, und rechnen positiv $x$ nach Süden, $y$ nach Westen, $z$ nach oben. Zugleich setzen wir
\[
\BareRoot{(xx+yy+zz)}=r,\quad x=r\cos f\cos g,\quad y=r\cos f\sin g,\quad z=r\sin f
\]
so dass $g$ das Azimuth der von der Mitte des Stabes nach dem fraglichen Punkte gezogenen geraden Linie, und $f$ ihre Neigung gegen die Horizontalebene bedeutet.

Wir bezeichnen ferner mit $M$ das absolute magnetische Moment des Magnetstabes; mit $F$ die Neigung seiner magnetischen Axe, positiv wenn der Nordpol höher liegt; mit $G$ das Azimuth dieser Axe. Zur Abkürzung schreiben wir
\[
\cos F\cos G=A,\qquad\cos F\sin G=B,\qquad\sin F=C
\]
wodurch also die magnetischen Momente des Magnetstabes relativ gegen die drei Coordinatenaxen beziehungsweise $MA,MB,MC$ werden.

Die von dem Magnetstabe in dem Punkte $x,y,z$ ausgeübte magnetische Kraft zerlegen wir parallel mit den drei Coordinatenaxen in die partiellen Kräfte $\xi,\eta,\zeta$.

Die ganze Intensität der reinen erdmagnetischen Kraft an diesem Orte bezeichnen wir mit $U$; ihren verticalen Theil mit $Z$; den horizontalen Theil $T$ zerlegen wir parallel mit den beiden ersten Coordinatenaxen in die partiellen Kräfte $X$ und $Y$. Alle diese Kräfte $\xi,\eta,\zeta,T,U,X,Y,Z$ sind homogene Grössen.

Endlich sei $i$ die magnetische Inclination, $D$ die Declination, wobei wir, um uns dem gewöhnlichen Gebrauche zu conformiren, $D$ von Norden nach Westen zählen, und $i$ wie positiv betrachten, wenn der Nordpol der Magnetnadel nach unten geneigt ist. Wir haben demnach
\[
\begin{aligned}
X&=-T\cos D=-U\cos i\cos D\\
Y&=T\sin D=U\cos i\sin D\\
Z&=-T\operatorname{tang}i=-U\sin i
\end{aligned}
\]

\Article{2}
Die Wirkung des Magnetstabes in dem Platze $x,y,z$ besteht in geringen Veränderungen der Bestimmungsstücke der erdmagnetischen Kraft, welche wir, da sie wegen ihrer Kleinheit unbedenklich nach den Regeln der Differentialrech%
\SourceLeaf{435}nung behandelt werden können, durch die vorgesetzte Charakteristik $\mathrm{d}$ bezeichnen wollen. Da nun
\[
\mathrm{d}X=\xi,\qquad\mathrm{d}Y=\eta,\qquad\mathrm{d}Z=\zeta
\]
so wird
\[
\begin{aligned}
\xi&=T\sin D.\mathrm{d}D-\cos D.\mathrm{d}T\\
\eta&=T\cos D.\mathrm{d}D+\sin D.\mathrm{d}T\\
\zeta&=-T\sec i^2\mathrm{d}i-\operatorname{tg}i.\mathrm{d}T
\end{aligned}
\]
woraus
\[
\begin{aligned}
\mathrm{d}D&=\frac{\sin D}{T}.\xi+\frac{\cos D}{T}.\eta\\
\mathrm{d}T&=-\cos D.\xi+\sin D.\eta\\
\mathrm{d}i&=-\frac{\cos i^2}{T}.\zeta-\frac{\sin2i}{2T}.\mathrm{d}T
\end{aligned}
\]
Endlich wird
\[
\mathrm{d}U=\cos i.\mathrm{d}T-\sin i.\zeta
\]
oder
\[
\frac{\mathrm{d}U}{U}=\frac{\cos i^2}{T}.\mathrm{d}T-\frac{\sin2i}{2T}.\zeta=\frac{\mathrm{d}T}{T}+\operatorname{tang}i.\mathrm{d}i
\]

\Article{3}
Das Potential der in dem Magnetstabe enthaltenen magnetischen Flüssigkeiten, in dem Punkte $x,y,z$, lässt sich in eine nach den Potenzen von $\frac1r$ fortschreitende Reihe entwickeln, von welcher für unsern Zweck bloss das Hauptglied beibehalten zu werden braucht, welches von der Ordnung $\frac{1}{rr}$ ist. Bezeichnen wir dies Potential mit $V$, so sieht man leicht, dass unter dieser Einschränkung
\[
V=\frac{M(Ax+By+Cz)}{r^3}
\]
wird. Bekanntlich erhält man $\xi,\eta,\zeta$ durch die partiellen Differentialquotienten von $V$ nach $x,y,z$; es ist nemlich
\[
\xi=-\frac{\mathrm{d}V}{\mathrm{d}x},\qquad\eta=-\frac{\mathrm{d}V}{\mathrm{d}y},\qquad\zeta=-\frac{\mathrm{d}V}{\mathrm{d}z}
\]
\SourceLeaf{436}folglich, wenn man, um abzukürzen,
\[
\frac{Ax+By+Cz}{r}=k
\]
setzt, und erwägt, dass die partiellen Differentialquotienten $\frac{\mathrm{d}r}{\mathrm{d}x},\frac{\mathrm{d}r}{\mathrm{d}y},\frac{\mathrm{d}r}{\mathrm{d}z}$ beziehungsweise $=\frac{x}{r},\frac{y}{r},\frac{z}{r}$ sind,
\[
\begin{aligned}
\xi&=\frac{3Mkx}{r^4}-\frac{MA}{r^3}\\
\eta&=\frac{3Mky}{r^4}-\frac{MB}{r^3}\\
\zeta&=\frac{3Mkz}{r^4}-\frac{MC}{r^3}
\end{aligned}
\]
Substituirt man also für $x,y,z,A,B,C$ ihre Werthe, so erhält man
\[
\begin{aligned}
k&=\cos f\cos F\cos(G-g)+\sin f\sin F\\
\mathrm{d}D&=\frac{M}{Tr^3}(3k\cos f\sin(D+g)-\cos F\sin(D+G))\\
\frac{\mathrm{d}T}{T}&=\frac{M}{Tr^3}(-3k\cos f\cos(D+g)+\cos F\cos(D+G))\\
\mathrm{d}i&=-\tfrac12\sin2i.\frac{\mathrm{d}T}{T}-\frac{M}{Tr^3}.\cos i^2(3k\sin f-\sin F)\\
\frac{\mathrm{d}U}{U}&=\cos i^2.\frac{\mathrm{d}T}{T}-\frac{M}{2Tr^3}.\sin2i(3k\sin f-\sin F)
\end{aligned}
\]
welche Formeln die vollständige Auflösung unsrer Aufgabe enthalten. Ohne unser Erinnern sieht man, dass $\mathrm{d}D$ und $\mathrm{d}i$ hier in Theilen des Halbmessers ausgedrückt sind, und also den Werthen noch der Factor $206265''$ beigefügt werden muss, um jene Änderungen in Bogentheilen zu erhalten. Der Werth von $\frac{M}{T}$ wird übrigens durch Versuche nach der in der \textit{Intensitas vis magneticae terrestris} gelehrten Methode bestimmt werden müssen.

\Article{4}
In der Ausübung sind solche Fälle die häufigsten, wo unsre allgemeinen Formeln durch specielle Verhältnisse eine bedeutende Vereinfachung erhalten. Es verdienen hier besonders die beiden folgenden bemerkt zu werden.

I. Wenn der Magnetstab vertical, also $F=\pm90^0$ ist, so nehmen die allgemeinen Formeln folgende Gestalt an:
\SourceLeaf{437}
\[
\begin{aligned}
k&=\pm\sin f\\
\mathrm{d}D&=\pm\frac{3M}{2Tr^3}\sin2f\sin(D+g)\\
\frac{\mathrm{d}T}{T}&=\mp\frac{3M}{2Tr^3}\sin2f\sin(D+g)\\
\mathrm{d}i&=-\tfrac12\sin2i.\frac{\mathrm{d}T}{T}\mp\frac{M}{Tr^3}\cos i^2(3\sin f^2-1)\\
\frac{\mathrm{d}U}{U}&=\cos i^2.\frac{\mathrm{d}T}{T}\mp\frac{M}{2Tr^3}\sin2i(3\sin f^2-1)
\end{aligned}
\]
Liegt zugleich der Punkt $x,y,z$ mit der Mitte des Magnetstabes in gleicher Höhe, so wird $z=0,f=0$ und folglich
\[
\begin{aligned}
\mathrm{d}D&=0\\
\mathrm{d}T&=0\\
\mathrm{d}i&=\pm\frac{M}{Tr^3}\cos i^2\\
\frac{\mathrm{d}U}{U}&=\pm\frac{M}{2Tr^3}\sin2i
\end{aligned}
\]
Es erhellt daraus, dass die Beobachtungen an einem Unifilar- oder Bifilarmagnetometer durch einen in demselben Locale befindlichen zweiten Magnetstab gar nicht gestört werden, so lange derselbe in verticaler Stellung und seine Mitte in derselben Höhe mit dem Stabe des Magnetometers erhalten wird.

II. Ist der Magnetstab horizontal, oder $F=0$, so gehen unsre Formeln in folgende über:
\[
\begin{aligned}
k&=\cos f\cos(G-g)\\
\mathrm{d}D&=\frac{M}{Tr^3}(3\cos f^2\cos(G-g)\sin(D+g)-\sin(D+G))\\
\frac{\mathrm{d}T}{T}&=\frac{M}{Tr^3}(-3\cos f^2\cos(G-g)\cos(D+g)+\cos(D+G))\\
\mathrm{d}i&=-\tfrac12\sin2i.\frac{\mathrm{d}T}{T}-\frac{3M}{2Tr^3}\cos i^2\sin2f\cos(G-g)\\
\frac{\mathrm{d}U}{U}&=\cos i^2.\frac{\mathrm{d}T}{T}-\frac{3M}{4Tr^3}\sin2i\sin2f\cos(G-g)
\end{aligned}
\]
Ist zugleich der Magnetstab im magnetischen Meridian (also $G=180^0-D$ für die natürliche Lage), oder senkrecht gegen denselben (also $G=90^0-D$ oder $=270^0-D$, jenachdem der Nordpol auf der Westseite oder auf der Ostseite sich befindet), so erhalten offenbar die Formeln für $\mathrm{d}D$ und $\mathrm{d}T$ noch weitere Ver%
\SourceLeaf{438}einfachung; diese Fälle treten ein, wenn der Stab, dessen Wirkung in der Ferne gesucht wird, den Bestandtheil eines Unifilar- oder eines Bifilarmagnetometers in transversaler Stellung ausmacht.

\Article{5}
Wenn man die Wirkung eines Magnetstabes in verschiedenen horizontalen Lagen unter einander vergleichen will, so kann man jeder der im vorhergehenden Artikel II gegebenen Formeln leicht eine dazu zweckmässige Gestalt geben. Bestimmt man z. B. zwei Grössen $p,P$ durch die Gleichungen
\[
\begin{aligned}
p\cos P&=(3\cos f^2-1)\sin(D+g)\\
p\sin P&=\cos(D+g)
\end{aligned}
\]
so verwandelt sich die Formel für $\mathrm{d}D$ in folgende
\[
\mathrm{d}D=\frac{Mp}{Tr^3}\cos(G-g+P)
\]
woraus erhellt, dass $\mathrm{d}D$ für $G=g-P$ oder für $G=180^0+g-P$ seinen grössten Werth $\frac{Mp}{Tr^3}$ mit positivem oder negativem Zeichen erhält, hingegen für $G=90^0+g-P$ und für $G=270^0+g-P$ verschwindet. Auf gleiche Weise wird, wenn man
\[
\begin{aligned}
q\cos Q&=(3\cos f^2-1)\cos(D+g)\\
q\sin Q&=\sin(D+g)
\end{aligned}
\]
setzt,
\[
\frac{\mathrm{d}T}{T}=-\frac{Mq}{Tr^3}\cos(G-g-Q)
\]
woraus für den Maximumwerth und das Verschwinden ähnliche Bestimmungen hervorgehen.

\Article{6}
Die hiesigen Einrichtungen bieten zu einer mehrfachen Anwendung der gegebenen Vorschriften Gelegenheit dar, bei Bestimmung der wechselseitigen Einwirkung der Magnetstäbe des Unifilar- und des Bifilarmagnetometers auf einander, und der Wirkungen beider Stäbe an einem dritten Platze, wo auf einem festen Steinpostamente mit andern Apparaten von Zeit zu Zeit magnetische Beobachtungen im Freien gemacht werden. Die in Metern ausgedrückten auf die Mitte
\SourceLeaf{439}der Axe des \textsc{Reichenbach}schen Meridiankreises, und rücksichtlich der dritten Coordinate auf den Fussboden der Sternwarte bezogenen absoluten Coordinaten dieser drei Plätze sind folgende:

(I) Mitte des fünfundzwanzigpfündigen Magnetstabes des Bifilarmagnetometers, für welchen, das Meter als Längeneinheit angenommen, $\frac{M}{T}=2.63318$ ist,
\[
x=-3.391,\qquad y=+6.708,\qquad z=+0.661
\]
(II) Mitte des vierpfündigen Magnetstabes des Unifilarmagnetometers, für welchen $\frac{M}{T}=0.48592$
\[
x=-23.618,\qquad y=+69.206,\qquad z=-2.235
\]
(III) Mitte des Steinpostaments, und rücksichtlich der Höhe, Platz welchen die Mitte der Nadel eines \textsc{Robinson}schen Inclinatoriums einnimmt,
\[
x=-21.546,\qquad y=+56.979,\qquad z=-1.665
\]
Hier mögen nur die Endresultate einer vierfachen Rechnung Platz finden, in welcher für $D$ und $i$ die Werthe $18^0 11'54''$ und $67^0 36'$ zum Grunde gelegt sind. Die Veränderlichkeit dieser Grössen, so wie der Werthe von $\frac{M}{T}$ für die beiden Magnetstäbe kommt für den gegenwärtigen Zweck nicht in Betracht.

(1) und (4) Wirkungen des Magnetstabes des Unifilarmagnetometers, jene an dem Platze (III), diese an dem Platze (I).

(2) und (3) Wirkungen des Magnetstabes des Bifilarmagnetometers an den Plätzen (III) und (II), indem jener Stab in der transversalen Lage, Nordpol im Westen vorausgesetzt wird.
\[
\begin{array}{c|r|r|r|r}
&\mathrm{d}D&\mathrm{d}T&\mathrm{d}i&\mathrm{d}U\\\hline
(1)&+64''72&-0.0000884\,T&+6''91&-0.0000071\,U\\
(2)&+3.04&+0.0000250\,T&-1.76&+0.0000043\,U\\
(3)&+1.82&+0.0000132\,T&-0.93&+0.0000023\,U\\
(4)&+0.50&+0.0000001\,T&-0.00&+0.0000001\,U
\end{array}
\]
Die Zahlen für (2) und (3) verändern bloss ihre Zeichen, wenn im Bifilarmagnetometer der Stab die transversale Lage Nordpol Ost hat. Es beträgt also die ganze Störung an dem Platze III durch beide Apparate
\SourceLeaf{440}
\[
\begin{array}{c|r|r}
&\multicolumn{2}{c}{\text{Nordpol im Bifilarmagnetometer.}}\\
&\text{West}&\text{Ost}\\\hline
\mathrm{d}D&+67''76&+61''68\\
\mathrm{d}T&-0.0000634\,T&-0.0001134\,T\\
\mathrm{d}i&+5''15&+8''67\\
\mathrm{d}U&-0.0000028\,U&-0.0000114\,U
\end{array}
\]

\Article{7}
Schliesslich soll hier noch der Zusammenhang der im 2. Art. für die Wirkung eines Magnetstabes in der Ferne gegebenen Formeln mit einer einfachen schon im 2. Bande der \textit{Resultate} [1837. II.] erwähnten Construction gezeigt werden. Eine Figur kann man entweder nach den folgenden Angaben sich leicht selbst entwerfen, oder a. a. O. nachsehen.

Es sei $A$ der Mittelpunkt des Magnetstabes, $n$ ein beliebiger anderer Punkt seiner durch $A$ gelegten magnetischen Axe auf der Seite des magnetischen Nordpols, $s$ ein ähnlicher Punkt auf der Seite des Südpols, $C$ der Punkt, für welchen die magnetische Wirkung des Magnetstabs auf die daselbst concentrirt gedachte Einheit des nördlichen magnetischen Fluidums bestimmt werden soll. Die partiellen Kräfte $\xi,\eta,\zeta$ werden nach Art. 2 durch die Formeln ausgedrückt
\[
\begin{aligned}
\xi&=\frac{3Mkx}{r^4}-\frac{MA}{r^3}\\
\eta&=\frac{3Mky}{r^4}-\frac{MB}{r^3}\\
\zeta&=\frac{3Mkz}{r^4}-\frac{MC}{r^3}
\end{aligned}
\]
wo, wie man leicht sieht, $k$ dem Cosinus des Winkels zwischen $An$ und $AC$ gleich ist. Die ersten Theile von $\xi,\eta,\zeta$ vereinigen sich offenbar in Eine Kraft $\frac{3Mk}{r^3}$, die abstossend in der Richtung der geraden Linie $AC$ wirkt, wenn $k$ positiv ist, anziehend oder in der entgegengesetzten Richtung $CA$, wenn $k$ negativ ist. Eben so werden die zweiten Theile von $\xi,\eta,\zeta$ zu Einer Kraft $\frac{M}{r^3}$ zusammengesetzt, deren Richtung immer mit $ns$ parallel ist. Für den speciellen Fall, wo $AC$ mit der magnetischen Axe einen rechten Winkel macht, also $k=0$ ist, verschwindet die erste Kraft, und die zweite allein stellt also die
\SourceLeaf{441}ganze Wirkung dar. In jedem andern Falle sei in der Ebene, in welcher $n,A,s,C$ liegen, $CB$ eine Normale gegen $CA$, $B$ ihr Durchschnittspunkt mit der magnetischen Axe, und [$D$ derjenige Punkt auf der Geraden $AB$, für welchen] $AD=\tfrac13AB$. Für den Fall der Figur im 2. Bande der \textit{Resultate}, wo $AC$ mit $An$ einen stumpfen Winkel macht, also $D$ und $B$ auf der Seite des Südpols liegen, sind die beiden oben angegebenen Kräfte den geraden Linien $CA$ und $AD$ offenbar proportional, und der Richtung nach die erste mit $CA$ zusammenfallend, die andere mit $AD$ parallel; die Richtung ihrer Resultante wird also $CD$ und die Stärke derselben $=\frac{CD}{AD}.\frac{M}{r^3}$ sein. Für den andern in der Figur a. a. O. nicht gezeichneten Fall, wo $AC$ mit $An$ einen spitzen Winkel macht, also $B$ und $D$ auf der Seite des Nordpols liegen, findet dasselbe Resultat bloss mit dem Unterschiede Statt, dass die Richtung des Winkels des Magnetstabes auf ein Element nördlichen Fluidums nicht durch $CD$, sondern durch $DC$ ausgedrückt wird, was mithin a. a. O. zur Vervollständigung noch hinzugefügt werden muss.
